\documentclass[11pt]{article}
\usepackage[a4paper, margin=1in]{geometry}
\usepackage[T1]{fontenc}
\usepackage[utf8]{inputenc}
\usepackage{amsmath, amssymb}
\usepackage{hyperref}
\usepackage{xcolor}
\hypersetup{colorlinks=true, linkcolor=blue, urlcolor=blue, citecolor=blue}
\usepackage{authblk}
\usepackage{microtype}
\usepackage{enumitem}
\setlist[itemize]{leftmargin=1.4em}

\title{\textbf{Where $E=mc^2$ Comes From}\\[0.3em]
\large A foundations note: a bridge from the wave/radiation sector, not from RealQM or Newton}
\author[1]{Claes Johnson}
\affil[1]{KTH Royal Institute of Technology, Stockholm, Sweden \\ \texttt{claesjohnson@gmail.com}}
\date{\today \\[0.4em]\small\textit{Developed in collaboration with Claude (Anthropic).}}

\begin{document}
\maketitle

\begin{abstract}
$E=mc^2$ is \emph{foreign to both RealQM and Newton} --- it must be imported. It enters from the
\emph{wave/radiation} side (matter and light as waves), it is the \emph{one bridge} between the electric-energy
sector (RealQM / RealNucleus / Radiation) and the gravitational-mass sector, its factor $c^2$ is a \emph{unit
convention}, and the single non-conventional fact is that the exchange rate \emph{is the speed of light} ---
because matter and light are the same kind of wave.
\end{abstract}

\section{It is foreign to RealQM and to Newton}
RealQM is the \emph{static, non-relativistic} Coulomb structure: energies in Hartree (Coulomb attraction plus
the kinetic term $\kappa|\nabla\psi|^2$), with \textbf{no $c$ and no rest energy}; nothing in it produces
$mc^2$. Newtonian mechanics has only $\tfrac12 mv^2$ and potential energy --- \textbf{no rest energy at all}.
So $E=mc^2$ is not derivable within either; it is genuinely \textbf{added from outside} the matter-structure
sector.

\section{Where it does come from}
Not from the matter-structure side --- from the \textbf{wave/radiation} side, where $c$ lives.
\begin{itemize}
\item \textbf{Standard route:} special relativity (Lorentz invariance) --- which this program declines to use
as the derivation.
\item \textbf{In this program (Many-Minds Relativity):} from the wave picture --- but this route is itself
largely \emph{definitional}. The \textbf{physical input is $P=E/c$}, the momentum of a wave (radiation
pressure); the step $m:=P/c$ (i.e.\ $P=mc$) is a \textbf{definition of mass} --- a convention, not a theorem;
and then $E=Pc=mc^2$ follows algebraically. So the MMR ``derivation'' is \textbf{one physical fact ($P=E/c$)
plus one definitional convention ($m=P/c$)}.
\end{itemize}
So $E=mc^2$ appears exactly when the static Coulomb structure is made \textbf{dynamic and radiating}: it is a
statement about the \textbf{Radiation sector}, not the RealQM/atom sector --- and its derivation is, as just
noted, mostly convention.

\section{The content, though, is empirical: energy really is mass}
Convention governs the \emph{derivation} and the \emph{units}; the physical claim --- that all internal energy
contributes to mass --- is \textbf{empirical and confirmed}:
\begin{itemize}
\item \textbf{The proton is the proof.} Its mass is ${\sim}99\%$ the internal kinetic + binding energy of its
quarks and gluons (only ${\sim}1\%$ quark rest mass). A system's mass \emph{literally is} its internal energy
--- $E=mc^2$ read backwards, and measured.
\item \textbf{Nuclear binding is a weighed mass defect.} A deuteron weighs \emph{less} than a free proton +
neutron by $2.224$ MeV/$c^2$ --- bind (cool) the system and it weighs less, directly on the scale.
\item \textbf{So a cooling body does lose mass and become easier to accelerate.} Hot $=$ more internal energy
$=$ more mass $=$ more inertia $=$ harder to accelerate; cooling $\to$ lighter $\to$ easier. Real --- but for
ordinary \emph{thermal} energy, tiny: heating $1$ kg of water by $100$ K adds $\Delta E\approx4\times10^5$ J
$\to \Delta m\approx 5\times10^{-12}$ kg (${\sim}5$ ng), ${\sim}10^{-12}$ of the mass, far below direct
weighing. The \emph{confirmed} cases are nuclear (mass defect) and the proton.
\end{itemize}

\section{Is it ``only a convention''? --- half yes}
\begin{itemize}
\item \textbf{The $c^2$ factor is a unit convention.} In natural units ($c=1$), $E=m$: mass and energy are the
\emph{same quantity}, and $c^2$ is only the exchange rate that appears because mass (kg) and energy (J) are
historically measured in different units. The \emph{number} $c^2$ is a currency rate.
\item \textbf{The equivalence it encodes is not conventional.} That a bound system \emph{weighs less} (the mass
defect), that mass and radiation interconvert (annihilation, nuclear energy) --- that is empirical.
\item \textbf{In this program it is the postulated bridge.} $E=mc^2$ is the coupling \emph{axiom} tying the
electric-energy sector to the gravitational-mass sector --- not derived from either. It plays the role of a
\textbf{conversion rule in a conserved-energy ledger}, where energy runs \emph{electric (Coulomb)
$\leftrightarrow$ mass (gravitational) $\leftrightarrow$ radiation}, conserved overall, with $c^2$ the exchange
rate.
\end{itemize}
So ``a convention in a game of constant energy with two forms'' is a fair description of its \emph{function} ---
provided one keeps the physical equivalence (empirical) separate from the unit factor (conventional).

\section{The one non-conventional fact: the rate is the speed of light}
$c$ enters the program \textbf{only here} --- and it is the \textbf{same $c$} as light's wave speed. That
identity is the real content: the mass--energy exchange rate equals the speed of light because \textbf{matter
and light are both waves in the same arena, with one characteristic speed}. The units ($c^2$) are arbitrary;
that the rate is \emph{$c$ rather than some unrelated constant} is not --- it is the statement that matter and
radiation share one wave medium. Remove that shared-wave picture and $E=mc^2$ would be a bare convention; with
it, the $c$ is forced.

\section{Its role in the program}
$E=mc^2$ is the \textbf{seam where the Coulomb sector couples to gravitation}. Gravity never sees electric
charge directly; it sees only \textbf{energy}, and $E=mc^2$ turns the Coulomb sector's energy (binding,
kinetic, radiation) into \textbf{mass}, the gravitational charge. The coupling is \textbf{realized in matter}
(each cloud carries both a charge and a mass) and \textbf{manifested in the mass defect} (Coulomb binding
lowers the energy, hence the mass; the difference radiates). It is also the \textbf{only place $c$ enters} an
otherwise $c$-free, non-relativistic construction --- the relativistic seam of the whole program.

\medskip\noindent\emph{In one line.} $E=mc^2$ is not from RealQM or Newton but imported from the
wave/radiation side ($P=mc$); its $c^2$ is a unit convention and it functions as the conversion rule in a
conserved-energy ledger between the electric and gravitational forms --- the one non-conventional fact being
that the rate \emph{is} the speed of light, because matter and light are the same kind of wave.

\end{document}
